\section{The Finite-Variance Classification}
\label{sec:finite-variance-equivalence-and-global-rigidity}

We first realize row series on a common measurable carrier and prove the
equivalent--singular orbit dichotomy.  Hilbert--Schmidt sufficiency follows by
finite-coordinate approximation.  For necessity, qualitative tail rigidity is
upgraded to square summability by conjugating remote coordinate swaps and
applying the local estimate from \cref{prop:local-fisher-hellinger}.

\subsection{Measurable row-series realizations on countable operator orbits}
\label{sec:measurable-realization}

Only \(\E X=0\) and \(\E X^2=1\) are needed here.  Each row series exists as
an \(L^2\) equivalence class.  The distinction between an operator on
\(\ell^2\) and its measurable realization on \(\R^\N\) is already central
to~\cite{shimomura1977linear}.  Such equivalence classes cannot be composed pointwise:
separately chosen versions for different operators or orbit laws may fail on
incompatible null sets.  We therefore select deterministic subsequences for
every row in a fixed countable operator group, so that operator multiplication
becomes composition on one common measurable carrier.  Countability is used
only to intersect the exceptional sets; no universal pointwise action of all of
\(\GL(H)\) is asserted.

\begin{lemma}[Measurable row-series realization on each countable operator orbit]
\label{lem:measurable-row-series-realization}
For every \(T\in\mathcal B(H)\), there is a measurable row-series
realization, namely a Borel map
\(\Phi_T:\R^{\N}\to\R^{\N}\) whose coordinates are the \(L^2(\mu)\) row sums
\[
  (\Phi_Tx)_i=\sum_{j\geq1}T_{ij}x_j.
\]
These maps can be chosen so that the following statements hold.
If a probability measure \(\lambda\) satisfies, for some finite
\(C_\lambda\geq0\),
\begin{equation}
  \int\left|\sum_j a_jx_j\right|^2\,d\lambda(x)
  \leq C_\lambda\|a\|_2^2
  \quad\text{for every finitely supported }a,
  \label{eq:covariance-domination}
\end{equation}
then \((\Phi_T)_\#\lambda\) also satisfies this condition.
Its constant is at most \(C_\lambda\|T\|_{\mathrm{op}}^2\).  Also
\begin{equation}
  \Phi_S\circ\Phi_T=\Phi_{ST}
  \quad\lambda\text{-almost surely}
  \qquad(S,T\in\mathcal B(H)).
  \label{eq:measurable-composition}
\end{equation}
For each fixed countable subgroup \(\Gamma\leq\GL(H)\),
there is a \(\Gamma\)-invariant Borel set \(D\subset\R^{\N}\) of full
measure for every \(\lambda\) satisfying
\eqref{eq:covariance-domination}.
On this set, these restrictions form an action by bimeasurable
automorphisms.
In particular, this applies to \(\mu\) and all its \(\Gamma\)-orbit laws,
including singular orbit laws.
\end{lemma}

\begin{proof}
For \(a\in\ell^2\), the finite sums \(\sum_{j\leq n}a_jx_j\) have a
unique \(L^2(\lambda)\) limit \(L_a\).  Choose
\[
  n_k(a)=\min\left\{n\geq k:
                \sum_{j>n}|a_j|^2\leq2^{-4k}\right\}.
\]
Define \(\widehat L_a(x)\) to be the limit of
\(\sum_{j\leq n_k(a)}a_jx_j\) when it exists, and zero otherwise.
By \eqref{eq:covariance-domination},
\[
  \left\|\sum_{j\leq n_k(a)}a_jx_j-L_a\right\|_{L^2(\lambda)}^2
  \leq C_\lambda2^{-4k}
\]
\(\lambda\)-almost surely.  This deterministic definition works for every
law satisfying \eqref{eq:covariance-domination}.  Set
\(\Phi_T(x)_i=\widehat L_{T^*e_i}(x)\); the construction is Borel.

For finitely supported \(b\),
\(\sum_i b_i\Phi_T(x)_i=L_{T^*b}\) almost surely, proving the asserted
second-moment bound for the pushforward.  For the \(i\)-th row of \(S\),
\[
  \sum_{j\leq n}S_{ij}\widehat L_{T^*e_j}
     =L_{T^*P_nS^*e_i}
     \longrightarrow L_{T^*S^*e_i}
     \quad\text{in }L^2(\lambda).
\]
The selected subsequence on the left also converges almost surely to
\(\Phi_S(\Phi_T(x))_i\), because the output law satisfies the same type of
bound.  Uniqueness of the limit proves
\eqref{eq:measurable-composition}.

For a countable \(\Gamma\), let \(B\) be the Borel set on which all its
row limits, identity relations, and pairwise composition relations hold.
The preceding argument gives \(\lambda(B)=1\) for every admissible law.  Put
\[
  D=\bigcap_{T\in\Gamma}\Phi_T^{-1}(B).
\]
Every orbit law satisfies \eqref{eq:covariance-domination}, so \(D\) has full
measure.  The composition identities make \(D\) invariant and
\(\Phi_{T^{-1}}\) the inverse of \(\Phi_T\) on \(D\).
\end{proof}

Write \(\nu_T=(\Phi_T)_\#\mu\).  Every countable operator family used below
therefore acts on one common measurable carrier.  No universal pointwise action of all
\(\GL(H)\) is needed; on the original product law these versions agree with
the usual almost-sure row sums.

\subsection{Rational translation invariance and the zero--one dichotomy}
\label{sec:translation-dichotomy}

Quasi-invariance under \(\ell^2\)-translations eliminates the intermediate
case: every bounded invertible orbit is either equivalent to \(\mu\) or
singular to it.
The translation input belongs to the Shepp-space problem of identifying
all shifts that preserve a product measure class
\cite{shepp1965translate}.  Beyond scalar translations, this question
extends to infinite products of commutative Lie groups and to
coordinatewise group actions~\cite{hora1991quasi,sadasue2008quasi}.
Within the translation problem, finer questions concern quasi-invariance
in all \(\ell^q\) directions, \(1\le q<2\)
\cite{arutyunyan2015spaces}, and the linearity and sequence-space
description of the Shepp space~\cite{honda2015approximations}.
Here its \(\ell^2\) quasi-invariance is combined with translation ergodicity;
the following proof specializes the equivalent--singular mechanism
of~\cite{okazaki1985dichotomy} to each bounded invertible orbit.

\begin{lemma}[Dichotomy for every bounded invertible orbit]
\label{lem:orbit-dichotomy}
For every \(T\in\GL(H)\),
\begin{equation}
  \nu_T\sim\mu\quad\text{or}\quad\nu_T\perp\mu.
  \label{eq:orbit-dichotomy}
\end{equation}
This conclusion uses only the a.e. positivity of \(f\), its centered finite second
moment, and \(\sqrt f\in W^{1,2}(\R)\).
In particular, it holds under \eqref{eq:main-assumptions}.
\end{lemma}

\begin{proof}
Put \(g=\sqrt f\).  The translation estimate
\[
  1-\Aff(f,f(\cdot-a))
     =\tfrac12\|g(\cdot-a)-g\|_2^2
     \leq\tfrac12\|g'\|_2^2a^2
\]
and Kakutani's theorem imply
\((x\mapsto x+a)_\#\mu\sim\mu\) for every \(a\in\ell^2\).
Since the scalar pairing of a row of \(T\) with \(T^{-1}a\) is absolutely
convergent,
\[
  \Phi_T(x+T^{-1}a)=\Phi_T(x)+a
  \quad\mu\text{-almost surely};
\]
hence \(\nu_T\) is also quasi-invariant under every \(\ell^2\) translation.

Let \(\mathcal A\) be the countable additive group generated by
\(\{qe_j,qTe_j:q\in\mathbb Q,\ j\geq1\}\).
Both \(\mu\) and \(\nu_T\) are quasi-invariant and ergodic under
translation by \(\mathcal A\).  Indeed, rational-translation invariance and
positivity of the finite-dimensional densities make every finite-coordinate
section constant almost everywhere (convolve with a smooth approximate
identity); the event is therefore a tail event, so Kolmogorov's zero--one law
applies.  Pullback by \(\Phi_T\) proves the same assertion for \(\nu_T\),
because translation by \(qTe_j\) pulls back to translation by \(qe_j\).

Write \(\nu_T=\nu_a+\nu_s\) for its Lebesgue decomposition relative to
\(\mu\).  Choose a Borel \(\mu\)-null set \(N_0\) supporting
\(\nu_s\), and let \(N=\bigcup_{a\in\mathcal A}(N_0+a)\).
Quasi-invariance of \(\mu\) makes \(N\) a \(\mu\)-null set;
it is exactly \(\mathcal A\)-invariant and still supports \(\nu_s\).
Thus \(\nu_T(N)=\nu_s(\R^{\N})\in\{0,1\}\) by ergodicity.  Mass one gives
singularity.  If the mass is zero, write \(\nu_T=\ell_T\mu\).
Quasi-invariance makes \(\{\ell_T>0\}\) invariant modulo \(\mu\); it has
positive measure and hence, by ergodicity, full measure.  Thus
\(\nu_T\sim\mu\).
\end{proof}

\subsection{Finite-marginal likelihood ratios and projected Hellinger energy}
\label{sec:finite-marginal-likelihood-ratios}

Finite-marginal likelihood ratios form a martingale, but its limit is the
density of the absolutely continuous part of \(\nu_T\), not automatically a
density of the whole law.  Martingale convergence therefore identifies the
limiting affinity; the orbit dichotomy is the additional input that upgrades
positive limiting affinity from mere overlap to equivalence.  This is why
bounded projected Hellinger energy is an exact diagnostic in the present orbit
problem.
For background on likelihood ratios and Hellinger comparison,
see~\cite{lecam1986asymptotic}; the role of Lebesgue decomposition in
equivalence and singularity criteria is discussed
in~\cite{mandrekar2018lebesgue}.

The convolution argument below establishes finite-dimensional absolute
continuity.  Global equivalence after an independent symmetric additive
perturbation is a distinct problem studied in~\cite{kitada1989absolute},
under a second-derivative integrability condition on the coordinate density.
For the dependent image law here, the global step uses translation ergodicity.

\begin{lemma}[Projected Hellinger energy and finite-marginal likelihood ratios]
\label{lem:finite-marginal-likelihood-ratios}
Assume \(f>0\) almost everywhere, \(\sqrt f\in W^{1,2}(\R)\),
and \(\E X=0\), \(\E X^2=1\).
For \(T\in\GL(H)\), every finite output law
\((\pi_n)_\#\nu_T\) is equivalent to \(\mu_n\).
Writing \(\mathcal F_n=\sigma(x_1,\ldots,x_n)\), define
\begin{equation}
 \begin{aligned}
  \ell_{T,n}(x)=\frac{d[(\pi_n)_\#\nu_T]}{d\mu_n}(\pi_n x),
  &\qquad E_n(T)=-2\log\E_\mu\sqrt{\ell_{T,n}},\\
  \Aff\bigl(\mu_n,(\pi_n)_\#\nu_T\bigr)
  &=\E_\mu\sqrt{\ell_{T,n}}.
 \end{aligned}
  \label{eq:finite-marginal-likelihood-ratios}
\end{equation}
Then \((\ell_{T,n})_{n\geq1}\) is a nonnegative mean-one
\(\mu\)-martingale.
We have
\begin{equation}
  \Aff\bigl(\mu_n,(\pi_n)_\#\nu_T\bigr)\downarrow\Aff(\mu,\nu_T),
  \qquad
  \nu_T\sim\mu
    \Longleftrightarrow\sup_nE_n(T)<\infty.
  \label{eq:marginal-affinity-limit}
\end{equation}
Otherwise \(E_n(T)\uparrow\infty\).
Then \(\nu_T\perp\mu\).
If \(\nu_T\sim\mu\) with density
\(\ell_T=d\nu_T/d\mu\), then
\begin{align}
  \ell_{T,n}&=\E_\mu[\ell_T\mid\mathcal F_n]
       \longrightarrow \ell_T
       &&\text{in }L^1(\mu),
       \label{eq:forward-likelihood-ui}\\
  \ell_{T,n}^{-1}&=\E_{\nu_T}[\ell_T^{-1}\mid\mathcal F_n]
       \longrightarrow \ell_T^{-1}
       &&\text{in }L^1(\nu_T).
       \label{eq:reverse-likelihood-ui}
\end{align}
In particular, both likelihood-ratio martingales are uniformly integrable in
their respective reference measures.
\end{lemma}

\begin{proof}
The first \(n\) rows of \(T\) are linearly independent, since \(T^*\)
is injective, so some \(n\) input columns form an invertible minor \(A\).
Writing the first \(n\) outputs as \(AU+W\), with \(U\sim\mu_n\) independent
of \(W\), shows that their density is a mixture of translates of the
a.e.-positive density of \(AU\).  It is positive and finite almost
everywhere, proving finite-dimensional equivalence.

Consistency of projected laws proves the martingale assertion.  The
martingale limit is
\(\ell_{T,\infty}=d(\nu_T)_a/d\mu\), the density of the absolutely continuous
part of \(\nu_T\).  Indeed, Fatou on cylinder sets gives
\(\ell_{T,\infty}\mu\leq\nu_T\); conversely, if
\((\nu_T)_a=\ell_{\mathrm{ac}}\mu\), then
\(\E_\mu[\ell_{\mathrm{ac}}\mid\mathcal F_n]\leq\ell_{T,n}\), and martingale
convergence gives the reverse inequality.  Since
\(\E_\mu(\sqrt{\ell_{T,n}})^2=1\), the square roots are uniformly integrable
and converge in \(L^1\) to \(\sqrt{\ell_{T,\infty}}\).  Conditional Jensen
also makes their expectations decreasing; hence
\[
  \lim_n\Aff\bigl(\mu_n,(\pi_n)_\#\nu_T\bigr)
      =\int\sqrt{\ell_{T,\infty}}\,d\mu
      =\Aff(\mu,\nu_T).
\]
The dichotomy in \cref{lem:orbit-dichotomy} now proves
\eqref{eq:marginal-affinity-limit} and the singular alternative.

Under equivalence, the forward identity is the defining
conditional-expectation identity for a marginal density.  For the reverse,
for \(B\in\mathcal F_n\),
\[
  \int_B \ell_{T,n}^{-1}\,d\nu_T=\mu(B)
       =\int_B \ell_T^{-1}\,d\nu_T.
\]
Conditional-expectation convergence gives both \(L^1\) limits and hence
both uniform-integrability assertions.
\end{proof}

Because initial segments are cofinal among finite coordinate sets, data
processing also shows that the criterion is independent of the coordinate
enumeration.
Here data processing is the contraction of Hellinger divergence under
measurable projections, a special case of the divergence inequalities
in~\cite{liese2006divergences}.

\subsection{Sufficiency of Hilbert--Schmidt perturbations}
\label{sec:sufficiency}

Finite-coordinate Hilbert--Schmidt approximants are Cauchy in Hellinger
distance, but this alone produces only a candidate density and one direction
of absolute continuity.  The inverse approximants and the common measurable
carrier are what prevent the limit from losing mass in the reverse direction.

\begin{proposition}[Hilbert--Schmidt sufficiency and likelihood-ratio convergence]
\label{prop:hilbert-schmidt-sufficiency}
Assume \eqref{eq:fisher-regularity}, let
$\mu=\rho^{\otimes\N}$, where $\rho(dx)=f(x)\,dx$, and let
$T=I+K\in\GL(H)$ with $K\in\mathcal S_2$.
Then \(\nu_T\sim\mu\).
For the eventually invertible operators
$T_n=I+P_nKP_n$, their likelihood ratios
$\ell_{T_n}=d\nu_{T_n}/d\mu$ converge in $L^1(\mu)$ to the strictly
positive likelihood ratio $\ell_T=d\nu_T/d\mu$, and
$\sqrt{\ell_{T_n}}\to\sqrt{\ell_T}$ in $L^2(\mu)$.
The inverse approximants satisfy the same conclusion:
\[
 \ell_{T_n^{-1}}\longrightarrow\ell_{T^{-1}}
 \quad\text{in }L^1(\mu).
\]
For $S,T\in\GL(H)$ with $S-I,T-I\in\mathcal S_2$, the
likelihood ratios obey, on the common measurable carrier of their countable
operator group,
\begin{align}
 \ell_{ST}(y)&=\ell_S(y)\ell_T(\Phi_{S^{-1}}y),
   \label{eq:rn-cocycle}\\
 \ell_T(y)\ell_{T^{-1}}(\Phi_{T^{-1}}y)&=1.
   \label{eq:rn-inverse}
\end{align}
If $Q$ is an exact linear stabilizer of $\mu$, the same equivalence
holds whenever \(TQ^{-1}-I\in\mathcal S_2\) and \(T\) is invertible.
\end{proposition}

\begin{proof}
Since $P_nKP_n\to K$ in Hilbert--Schmidt norm, $T_n\to T$ in
operator norm, so the approximants are eventually invertible with uniformly
bounded inverse norms; their finite-coordinate likelihood ratios
\(\ell_{T_n}\) are positive almost everywhere.
For all sufficiently large $m,n$, the operator $T_n^{-1}T_m-I$
acts in a finite coordinate block and has operator norm at most $1/2$.
Hellinger invariance under the finite invertible map $T_n^{-1}$ and
\eqref{eq:finite-hellinger-upper} give
\begin{equation}
 \|\sqrt{\ell_{T_m}}-\sqrt{\ell_{T_n}}\|_2
 \le\sqrt{C_f}\,\|T_n^{-1}\|_{\mathrm{op}}
                    \|T_m-T_n\|_{\mathcal S_2}.
 \label{eq:approximant-hellinger-cauchy}
\end{equation}
Thus \(\sqrt{\ell_{T_n}}\) converges in $L^2(\mu)$ to a nonnegative
function of norm one; denote its square by $\widetilde\ell_T$.  Since
\[
 \|\ell_{T_m}-\ell_{T_n}\|_1
 \le2\|\sqrt{\ell_{T_m}}-\sqrt{\ell_{T_n}}\|_2,
\]
we also have \(\ell_{T_n}\to\widetilde\ell_T\) in \(L^1\).  Meanwhile,
\(\Phi_{T_n}\mathbf X\to\Phi_T\mathbf X\) in \(L^2\) on every finite set
of output coordinates.  Cylinder tests therefore identify
$\widetilde\ell_T\mu$ with $\nu_T$.

The inverse approximants converge in Hilbert--Schmidt norm because
\[
 \|T_n^{-1}-T^{-1}\|_{\mathcal S_2}
 \le\|T_n^{-1}\|_{\mathrm{op}}
       \|T_n-T\|_{\mathcal S_2}\|T^{-1}\|_{\mathrm{op}}.
\]
Applying the same argument gives $\nu_{T^{-1}}\ll\mu$.  On the common
measurable carrier supplied by
\cref{lem:measurable-row-series-realization},
\[
 \mu=(\Phi_T)_\#\nu_{T^{-1}}
       \ll(\Phi_T)_\#\mu=\nu_T.
\]
This proves equivalence and identifies
$\widetilde\ell_T=\ell_T>0$ almost everywhere.
Change of variables on that carrier proves
\eqref{eq:rn-cocycle} and \eqref{eq:rn-inverse}.  Finally, $\nu_Q=\mu$ and
$T=(TQ^{-1})Q$ give the assertion about exact stabilizers.
\end{proof}

\subsubsection{The orthogonal case requires only location Fisher information}
\label{sec:orthogonal-transformations}

For orthogonal transformations the assumptions can be reduced to
\begin{equation}
 f>0\ \text{almost everywhere},\qquad
 g=\sqrt f\in W^{1,2}(\R),\qquad
 \E X=0,\qquad\E X^2=1.
 \label{eq:orthogonal-fisher-regularity}
\end{equation}
In particular, $xg'\in L^2$ is not required.
The cutoff identities $\E s(X)=0$ and $\E[Xs(X)]=1$ still follow
from $J<\infty$ and the second moment.
For a skew-symmetric matrix $A$, the score has only off-diagonal terms and
\begin{equation}
 \E S_A^2=(J-1)\|A\|_F^2.
 \label{eq:orthogonal-fisher}
\end{equation}
The generator formula remains valid.
Only $x_j\partial_i\Omega_n$ with $i\ne j$ occur.
One can justify it by radial cutoff and radial mollification.
Those operators commute with the orthogonal one-parameter group.

\begin{proposition}[Orthogonal local bounds and Hilbert--Schmidt sufficiency]
\label{prop:orthogonal-sufficiency}
Under \eqref{eq:orthogonal-fisher-regularity}, every
orthogonal $U$ on $\ell^2$ with $U-I\in\mathcal S_2$ satisfies
$\nu_U\sim\mu$.  Its likelihood ratios have the same $L^1$ and
square-root $L^2$ convergence along finite-coordinate orthogonal
approximants, as do the inverse likelihood ratios.
If $f$ is non-Gaussian, there also exist dimension-free
$c,C,\delta>0$ such that
\begin{equation}
 c\|V-I\|_F\le h(V_\#\mu_n,\mu_n)
                   \le C\|V-I\|_F
 \quad(V\in O(n),\ \|V-I\|_F\le\delta).
 \label{eq:orthogonal-local-hellinger}
\end{equation}
\end{proposition}

\begin{proof}
For a skew-symmetric $A$, integration along $\exp(tA)$ gives
\begin{equation}
 h((\exp A)_\#\mu_n,\mu_n)
 \le\tfrac12\sqrt{J-1}\,\|A\|_F.
 \label{eq:orthogonal-path-upper}
\end{equation}
The construction of $p,q$ in \cref{lem:fourier-dual-moments} uses only
$s(X),X\in L^2$ and $J-1>0$.  For skew-symmetric $A$, take
$H_A=\sum_{i\ne j}A_{ij}p(X_i)q(X_j)$.  Both $H_A$ and $S_A$ have no
diagonal term, so the generator proof uses only bounded functions times
$s(X)$ or $X$ and remains valid under
\eqref{eq:orthogonal-fisher-regularity}.  The half-density pairing gives the
local lower bound in exponential coordinates.  Since a sufficiently small
orthogonal $V-I$ has a real skew-symmetric logarithm with
$\|\log V\|_F\asymp\|V-I\|_F$, this proves
\eqref{eq:orthogonal-local-hellinger}.

For the infinite-dimensional assertion, set $T_n=I+P_n(U-I)P_n$ and take
its polar factor
\[
 V_n=T_n(T_n^*T_n)^{-1/2}.
\]
It is orthogonal and acts identically outside the first $n$
coordinates.  Since $T_n\to U$ in Hilbert--Schmidt norm and its least
singular value is eventually bounded below by some $a>0$,
\[
 \|V_n-T_n\|_{\mathcal S_2}
 =\|I-|T_n|\|_{\mathcal S_2}
 \le\frac{\|T_n^*T_n-I\|_{\mathcal S_2}}{1+a}.
\]
The right side tends to zero, so $V_n\to U$ in Hilbert--Schmidt norm.
For large $m,n$, apply \eqref{eq:orthogonal-path-upper} to the small
skew-symmetric logarithm of $V_n^*V_m$.  The square roots of the likelihood
ratios are Cauchy, and cylinder tests identify their limit as in
\cref{prop:hilbert-schmidt-sufficiency}.  Repeating the argument for
$V_n^{-1}\to U^{-1}$ and using the measurable inverse action proves
equivalence and all asserted convergence.
\end{proof}

\subsection{Necessity of Hilbert--Schmidt perturbations}
\label{sec:necessity}

Approximation and likelihood-ratio convergence transfer local coercivity from
coordinate blocks to Hilbert--Schmidt operators, allowing many small tail
mismatches to be combined into a contradiction.

\begin{lemma}[The local bound for Hilbert--Schmidt operators]
\label{lem:local-hs-coercivity}
Let $f$ be non-Gaussian and satisfy
\eqref{eq:fisher-regularity}.  There are constants
$c,C,\delta>0$ such that every $K\in\mathcal S_2(\ell^2)$ with
$\|K\|_{\mathcal S_2}\le\delta$ satisfies
\begin{equation}
 c\|K\|_{\mathcal S_2}
 \le h(\nu_{I+K},\mu)
 \le C\|K\|_{\mathcal S_2}.
 \label{eq:local-hs-hellinger}
\end{equation}
The constants can be chosen with $\delta<1$, so $I+K$ is
automatically invertible.  Under the weaker assumptions
\eqref{eq:orthogonal-fisher-regularity}, the same conclusion
holds for orthogonal $I+K$ whenever $f$ is non-Gaussian.
This includes finite-rank perturbations whose ranges have
infinite coordinate support.
\end{lemma}

\begin{proof}
Let $K_n=P_nKP_n$.  The untouched tail has affinity one, so
\eqref{eq:local-matrix-hellinger} applies directly to
$h(\nu_{I+K_n},\mu)$.  By
\cref{prop:hilbert-schmidt-sufficiency},
\[
 h(\nu_{I+K_n},\mu)
 \longrightarrow h(\nu_{I+K},\mu),
 \qquad \|K_n\|_{\mathcal S_2}\longrightarrow\|K\|_{\mathcal S_2}.
\]
Passing to the limit proves both bounds.

In the orthogonal case, the polar approximants from
\cref{prop:orthogonal-sufficiency} satisfy
\[
 h(\nu_{V_n},\mu)\longrightarrow h(\nu_{I+K},\mu),
 \qquad
 \|V_n-I\|_{\mathcal S_2}\longrightarrow\|K\|_{\mathcal S_2}.
\]
After reducing the radius if necessary, pass to the limit in
\eqref{eq:orthogonal-local-hellinger}.
\end{proof}

\subsubsection{Characteristic-function rigidity and bounded invertible linear mixing}
\label{sec:characteristic-function-rigidity}

Necessity begins by identifying what a single output row can look like if
its law is close to the original coordinate law.
This is a one-dimensional rigidity problem for linear combinations of
independent copies of \(X\).
The characteristic-function argument below says that, for a non-Gaussian
coordinate law, such a row must asymptotically be one signed coordinate.

Throughout this subsection, $X_j$ are independent copies of a real random
variable $X$ with law $\rho$ and
\[
  \E X=0,\qquad \E X^2=1,
\]
and $\rho$ is non-Gaussian.  The characteristic-function arguments
below require no moments of order greater than two.  Recall
\[
  \phi(t)=\E e^{itX},\qquad \mu=\rho^{\otimes\N}.
\]
The group $\mathcal P$ consists of coordinate permutations whose row
signs preserve $\rho$.  This definition applies even when $\rho$ has
no density.

\begin{lemma}[A product approximation for small coefficients]
\label{lem:characteristic-product-tail}
There is a nonnegative function $\omega_\rho$, with
$\omega_\rho(\delta)\to0$ as
$\delta\downarrow0$, such that, for every finite or countable real family
$(b_j)$ with $\sum_j b_j^2<\infty$,
\begin{equation}
 \left|\prod_j\phi(tb_j)
       -\exp\!\left(-\frac{t^2}{2}\sum_jb_j^2\right)\right|
 \leq \omega_\rho(\delta)t^2\sum_jb_j^2
 \quad\text{if}\quad \sup_j|tb_j|\leq\delta.
 \label{eq:characteristic-product-tail}
\end{equation}
For a countable family, the product denotes the characteristic function
of the $L^2$ sum $\sum_jb_jX_j$.
\end{lemma}
\begin{proof}
The centered second-moment assumption gives
\[
 \phi(u)=1-\frac{u^2}{2}+o(u^2).
\]
Consequently
\[
 \omega_\rho(\delta):=
 \sup_{0<|u|\leq\delta}
 \frac{|\phi(u)-e^{-u^2/2}|}{u^2}
 \longrightarrow0,\qquad \omega_\rho(0):=0.
\]
Both factors have modulus at most one.  Telescoping finite products therefore
bounds their difference by
$\sum_j\omega_\rho(\delta)t^2b_j^2$, proving
\eqref{eq:characteristic-product-tail}.  The countable case follows from
$L^2$ convergence of the partial sums.
\end{proof}

\begin{lemma}[Qualitative rigidity of a unit row]
\label{lem:unit-row-rigidity}
Let $a^{(n)}\in\ell^2$ satisfy $\|a^{(n)}\|_2=1$.  If
\[
 \Law\!\left(\sum_j a_j^{(n)}X_j\right)
       \Longrightarrow\rho,
\]
then
\begin{equation}
 \inf_{\substack{j\geq1,\ \varsigma\in\{-1,1\}\\
                  \Law(\varsigma X)=\rho}}
       \|a^{(n)}-\varsigma e_j\|_2\longrightarrow0.
 \label{eq:unit-row-qualitative}
\end{equation}
\end{lemma}
\begin{proof}
Reorder each row by decreasing absolute value; this preserves the law of the
sum and gives $|a_j^{(n)}|\le j^{-1/2}$.  From any subsequence, extract a
further one with $a_j^{(n)}\to a_j$ for every $j$.  Fatou gives
\[
 \tau_a:=\sum_j a_j^2\leq1.
\]
For fixed $m,t$, \cref{lem:characteristic-product-tail} replaces the tail
characteristic function by
\[
 \exp\!\left[-\frac{t^2}{2}
        \left(1-\sum_{j\leq m}(a_j^{(n)})^2\right)\right]
\]
with error at most $t^2\omega_\rho(|t|/\sqrt{m+1})$, uniformly in $n$.
Letting first $n\to\infty$ and then $m\to\infty$ gives
\begin{equation}
 e^{-(1-\tau_a)t^2/2}\prod_j\phi(a_jt).
 \label{eq:unit-row-gaussian-summand}
\end{equation}
Thus the limit is $\sum_j a_jX_j$ plus an independent
$N(0,1-\tau_a)$ variable.

Under the convergence hypothesis this limiting law is the law of $X$.
Suppose that no $|a_j|$ is one.  If all $a_j$ vanish,
\eqref{eq:unit-row-gaussian-summand} already makes $X$ Gaussian.
Otherwise $m_a:=\max_j|a_j|$ exists and lies in $(0,1)$.  For
$\mathbf j=(j_1,\ldots,j_k)\in\N^k$, set
$b_{\mathbf j}:=a_{j_1}\cdots a_{j_k}$.
Iterating the characteristic-function identity
\eqref{eq:unit-row-gaussian-summand} $k$ times gives
\begin{equation}
 \phi(t)=e^{-(1-\tau_a^k)t^2/2}
       \prod_{\mathbf j\in\N^k}\phi(tb_{\mathbf j}),
 \label{eq:smoothing-iteration}
\end{equation}
where
\[
 \sum_{\mathbf j} b_{\mathbf j}^2=\tau_a^k,
 \qquad \sup_{\mathbf j}|b_{\mathbf j}|\leq m_a^k.
\]
The substitution is legitimate in $L^2$ at every finite depth.  Applying
\cref{lem:characteristic-product-tail} to
\eqref{eq:smoothing-iteration} yields
\[
 |\phi(t)-e^{-t^2/2}|
 \leq t^2\tau_a^k\omega_\rho(|t|m_a^k)\longrightarrow0.
\]
This contradicts non-Gaussianity.  Hence some
$a_j=\varsigma\in\{-1,1\}$.  All other coefficients and the Gaussian
summand vanish, and the limiting identity gives $\Law(\varsigma X)=\rho$.

If \eqref{eq:unit-row-qualitative} failed, a subsequence would stay a fixed
distance from every allowed signed unit vector.  The preceding extraction
produces an allowed $\varsigma e_j$ with
\[
 \|a^{(n)}-\varsigma e_j\|_2^2
       =2-2\varsigma a_j^{(n)}\longrightarrow0.
\]
Undoing the permutations is a contradiction.
\end{proof}

\begin{lemma}[Uniform invariance under tail symmetries]
\label{lem:tail-symmetry-uniform}
Let $\mu$ be the i.i.d. product law and let $\nu\ll\mu$.  Let
$\mathcal T_N$ consist of the members of $\mathcal P$ that change only
finitely many coordinates and fix coordinates $1,\ldots,N$.
Then
\begin{equation}
 \sup_{S\in\mathcal T_N}h(\nu,S_\#\nu)\longrightarrow0.
 \label{eq:tail-symmetry-uniform}
\end{equation}
Also, the one-coordinate marginals of $\nu$ converge to the common
coordinate law of $\mu$, uniformly over coordinates greater than $N$.
\end{lemma}
\begin{proof}
Put $\ell=d\nu/d\mu$ and
$\ell_N=\E_\mu[\ell\mid\sigma(x_1,\ldots,x_N)]$.  The martingale
convergence theorem gives $\|\ell-\ell_N\|_1\to0$.  Every
$S\in\mathcal T_N$ preserves $\mu$ and satisfies
$\ell_N\circ S^{-1}=\ell_N$.  Thus
\[
 h(\nu,S_\#\nu)^2
 \leq\|\ell-\ell\circ S^{-1}\|_{L^1(\mu)}
 \leq2\|\ell-\ell_N\|_{L^1(\mu)},
\]
which proves \eqref{eq:tail-symmetry-uniform}.
For $i>N$, the marginal density is
$\E_\mu[\ell\mid x_i]$, while independence gives
$\E_\mu[\ell_N\mid x_i]=1$, so
\[
 \left\|\E_\mu[\ell\mid x_i]-1\right\|_{L^1(\mu)}
 \leq\|\ell-\ell_N\|_{L^1(\mu)}.
\]
Taking the supremum over $i>N$ proves the marginal assertion.
\end{proof}

\begin{lemma}[Compactness of the covariance defect]
\label{lem:compact-covariance-rigidity}
Let $T$ be bounded on $\ell^2$.  If $\nu_T\ll\mu$, then
\[
 TT^*-I\quad\text{is compact}.
\]
If $T$ is boundedly invertible, $T^*T-I$ is compact as well.
\end{lemma}
\begin{proof}
We first show that, for every $B<\infty$, the squares of
\begin{equation}
 Y_b:=\sum_jb_jX_j,\qquad \|b\|_2\leq B,
 \label{eq:uniform-linear-square-family}
\end{equation}
are uniformly integrable.  Let
\[
 X^{(M)}=X\mathbf1_{\{|X|\leq M\}}
           -\E[X\mathbf1_{\{|X|\leq M\}}],
 \qquad Y_b^{(M)}=\sum_jb_jX_j^{(M)}.
\]
The variables $X^{(M)}$ are bounded and centered, and converge to $X$
in $L^2$.  Independence gives
\[
 \sup_{\|b\|_2\leq B}\|Y_b-Y_b^{(M)}\|_2
 \leq B\|X-X^{(M)}\|_2\longrightarrow0.
\]
For fixed $M$, the fourth-moment formula for independent centered variables
gives
\[
 \sup_{\|b\|_2\leq B}\E|Y_b^{(M)}|^4<\infty.
\]
It also makes the finite partial sums Cauchy in $L^4$, so the bound holds for
countable sums.  Moreover,
\[
 \sup_{\|b\|_2\leq B}
 \bigl\|Y_b^2-(Y_b^{(M)})^2\bigr\|_1
 \leq B^2\|X-X^{(M)}\|_2
                  (\|X\|_2+\|X^{(M)}\|_2)
 \longrightarrow0.
\]
Thus $Y_b^2$ is uniformly approximated in $L^1$ by a uniformly integrable
family, proving the claim using only the second moment of $X$.

Write $\nu=\nu_T$, $\ell=d\nu/d\mu$, and
$\ell_N=\E_\mu[\ell\mid\sigma(x_1,\ldots,x_N)]$.
For a finitely supported unit vector $a$ supported on coordinates greater than
$N$, put
$F_a(x)=(\sum_i a_ix_i)^2$.  Under $\mu$, these are squares from
\eqref{eq:uniform-linear-square-family}; under $\nu$, they are the squares of
$Y_{T^*a}$ under $\mu$.  Since $\|T^*a\|_2\leq\|T\|_{\mathrm{op}}$, they are
uniformly integrable under both laws.  The bounded statistic $F_a\wedge L$
depends only on tail coordinates, so independence gives
\[
 \left|\E_\nu(F_a\wedge L)-\E_\mu(F_a\wedge L)\right|
 \leq L\|\ell-\ell_N\|_1.
\]
Letting first $N\to\infty$ and then $L\to\infty$, uniformly in $a$, gives
\[
 \sup_{\substack{\|a\|_2=1,\ a\text{ finitely supported}\\
                  \operatorname{supp}a\subset\{N+1,N+2,\ldots\}}}
 \left|\langle(TT^*-I)a,a\rangle\right|\longrightarrow0.
\]
Since $TT^*-I$ is self-adjoint, density of finite vectors in the tail
subspace implies
\[
 \|(I-P_N)(TT^*-I)(I-P_N)\|_{\mathrm{op}}\longrightarrow0.
\]
The complementary rows and columns are finite rank, so $TT^*-I$ is compact.
If $T$ is invertible, then
$T^*T-I=T^{-1}(TT^*-I)T$.
\end{proof}

We now pass from qualitative row matching to square-summable error.  Compactness
of the covariance defect normalizes the remote rows; tail marginal convergence and
\cref{lem:unit-row-rigidity} match each with an allowed signed coordinate;
bounded invertibility makes these matches eventually distinct; and
\(T^*T-I\) gives the corresponding column matching.  None of these facts
controls the sum of the errors.

Suppose that sum diverges.  Isolate a remote finite packet with fixed small
error energy and swap it with fresh, almost perfectly matched auxiliary
coordinates.  Conjugating the output swap through \(T\) and cancelling the
aligned input swap produces \(W=PT^{-1}ST\): a small Hilbert--Schmidt
perturbation of the identity whose squared norm is comparable to the packet's
error energy.
Local coercivity must detect \(W\), while absolute continuity makes the same
change asymptotically invisible because \(S\) is supported arbitrarily far in
the tail.  The block computation below has one purpose: to verify this norm
comparison without losing constants to the ambient dimension.

\begin{lemma}[Hilbert--Schmidt necessity under a second-moment hypothesis]
\label{lem:hilbert-schmidt-necessity}
Assume that the common law has an a.e.-positive density $f$, that
$g=\sqrt f\in W^{1,2}(\R)$, and that $xg'\in L^2(\R)$.
If $T\in\GL(H)$ and $\nu_T\sim\mu$, then
\begin{equation}
 T-Q\in\mathcal S_2
 \qquad\text{for some }Q\in\mathcal P.
 \label{eq:hilbert-schmidt-necessity}
\end{equation}
When $T$ is orthogonal, the condition $xg'\in L^2$ can be omitted.
\end{lemma}
\begin{proof}
By \cref{lem:compact-covariance-rigidity}, $\|T^*e_i\|_2\to1$, and by
\cref{lem:tail-symmetry-uniform}, the law of the corresponding row sum tends
to $\rho$.  Normalization changes that sum by $o_{L^2}(1)$, so
\cref{lem:unit-row-rigidity} gives indices $\iota(i)$ and $\rho$-preserving
signs $\varsigma_i$ such that
\begin{equation}
 \widehat e_i:=\varsigma_i e_{\iota(i)},\qquad
 \|T^*e_i-\widehat e_i\|_2^2\longrightarrow0.
 \label{eq:qualitative-tail-matching}
 \end{equation}
The indices $\iota(i)$ are eventually distinct: if
$\iota(i)=\iota(j)$, then
$\varsigma_iT^*e_i-\varsigma_jT^*e_j$ has arbitrarily small norm
by \eqref{eq:qualitative-tail-matching}, whereas bounded
invertibility gives the lower bound $\sqrt2/\|T^{-1}\|_{\mathrm{op}}$.
Fix $N_0$ so that $\iota$ is injective on $i>N_0$.

Since $\iota(i)$ leaves every finite set, compactness of $T^*T-I$ gives
\[
 \|T\widehat e_i\|_2^2\longrightarrow1.
\]
Together with $\langle T\widehat e_i,e_i\rangle
=\langle\widehat e_i,T^*e_i\rangle\to1$, this yields
\begin{equation}
 \operatorname{err}_i:=\|T^*e_i-\widehat e_i\|_2^2
       +\|T\widehat e_i-e_i\|_2^2\longrightarrow0.
 \label{eq:row-column-tail-errors}
\end{equation}

We claim that $\sum_{i>N_0}\operatorname{err}_i<\infty$.  Choose the local
constants in \cref{lem:local-hs-coercivity} so that
\begin{equation}
 h(\mu,\nu_W)\geq c_{\mathrm{HS}}\|W-I\|_{\mathcal S_2}
 \quad\text{when }W\in\GL(H),\quad
 \|W-I\|_{\mathcal S_2}\leq\delta_{\mathrm{HS}}.
 \label{eq:local-lower-used}
\end{equation}
Suppose the claim fails, and choose $\delta>0$ so that
$3\|T^{-1}\|_{\mathrm{op}}\delta<\delta_{\mathrm{HS}}$.
In an arbitrarily remote tail where $\operatorname{err}_i\le\delta^2$, choose
a finite set $E$ satisfying
\begin{equation}
 \delta^2\leq\sum_{i\in E}\operatorname{err}_i\leq2\delta^2.
 \label{eq:active-tail-error}
\end{equation}
Pair each $i\in E$ with a distinct later index $k_i$ such that, for
$G=\{k_i:i\in E\}$, the sum $\sum_{k\in G}\operatorname{err}_k$ is
arbitrarily small.

Let $S$ exchange each output pair $(e_i,e_{k_i})$, and let $P$ exchange
the aligned signed input pairs $(\widehat e_i,\widehat e_{k_i})$, acting
identically on their orthogonal complements.  Both are involutions in
$\mathcal P$, so they preserve $\mu$.  Define
\[
 W=P T^{-1} S T.
\]
Then
\begin{equation}
 W-I=P T^{-1}(ST-TP),\qquad
 \frac{\|ST-TP\|_{\mathcal S_2}}{\|T\|_{\mathrm{op}}}
 \leq\|W-I\|_{\mathcal S_2}
 \leq\|T^{-1}\|_{\mathrm{op}}\,\|ST-TP\|_{\mathcal S_2}.
 \label{eq:swap-relative-norm}
\end{equation}
Hellinger invariance and $P_\#\mu=\mu$ give
\begin{equation}
 h(\mu,\nu_W)=h(\nu_T,S_\#\nu_T).
 \label{eq:swap-hellinger-identity}
\end{equation}
All differences have finite rank, and
\cref{lem:measurable-row-series-realization} supplies the common carrier
needed for this identity.

To compare the norms, freeze the auxiliary matches by replacing in $T$ each
row and column indexed by $G$ with its matched coordinate vector:
\[
 T_0^*e_k=\widehat e_k,\qquad T_0\widehat e_k=e_k\qquad(k\in G).
\]
and retain the complementary block of $T$.  Orthogonal block decomposition
gives
\begin{equation}
 \|T-T_0\|_{\mathcal S_2}^2
 \leq\sum_{k\in G}
       \bigl(\|T^*e_k-\widehat e_k\|_2^2
                    +\|T\widehat e_k-e_k\|_2^2\bigr)
 =\sum_{k\in G}\operatorname{err}_k.
 \label{eq:auxiliary-coordinate-approximation}
\end{equation}

Using the decompositions $E\oplus(E\cup G)^c\oplus G$ on the output and the
corresponding aligned input coordinates, write
\[
 T_0=\begin{pmatrix}A&B&0\\ C&D&0\\0&0&I\end{pmatrix},
 \qquad
 ST_0-T_0P=
 \begin{pmatrix}
 0&-B&I-A\\
 C&0&-C\\
 A-I&B&0
 \end{pmatrix}.
\]
Here $B,C$ are Hilbert--Schmidt, and the error on $E$ is
\[
 \sum_{i\in E}
 \bigl(\|T_0^*e_i-\widehat e_i\|_2^2
                    +\|T_0\widehat e_i-e_i\|_2^2\bigr)
 =2\|A-I\|_F^2+\|B\|_{\mathcal S_2}^2+\|C\|_{\mathcal S_2}^2.
\]
while the displayed commutator gives
\[
 \|ST_0-T_0P\|_{\mathcal S_2}^2
 =2\bigl(\|A-I\|_F^2+\|B\|_{\mathcal S_2}^2
                         +\|C\|_{\mathcal S_2}^2\bigr),
\]
Therefore
\begin{equation}
 \begin{aligned}
 2\|A-I\|_F^2+\|B\|_{\mathcal S_2}^2+\|C\|_{\mathcal S_2}^2
 &\leq\|ST_0-T_0P\|_{\mathcal S_2}^2\\
 &\leq2\bigl(2\|A-I\|_F^2+\|B\|_{\mathcal S_2}^2
                              +\|C\|_{\mathcal S_2}^2\bigr).
 \end{aligned}
 \label{eq:swap-block-energy}
\end{equation}
As $\sum_{k\in G}\operatorname{err}_k\to0$, the first displayed error tends
to $\sum_{i\in E}\operatorname{err}_i$, while
\[
 \|(ST-TP)-(ST_0-T_0P)\|_{\mathcal S_2}
 \leq2\|T-T_0\|_{\mathcal S_2}\longrightarrow0.
\]
Hence \eqref{eq:active-tail-error} and \eqref{eq:swap-block-energy} allow the auxiliary
indices to be chosen so that
\[
 \frac{\delta}{2}\leq\|ST-TP\|_{\mathcal S_2}\leq3\delta.
\]
It follows from \eqref{eq:swap-relative-norm} that
\[
 \frac{\delta}{2\|T\|_{\mathrm{op}}}
 \leq\|W-I\|_{\mathcal S_2}
 \leq3\|T^{-1}\|_{\mathrm{op}}\delta<\delta_{\mathrm{HS}}.
\]
By \eqref{eq:local-lower-used},
$h(\mu,\nu_W)\geq c_{\mathrm{HS}}\delta/(2\|T\|_{\mathrm{op}})$.
But $S$ is supported in the chosen remote tail, so
\cref{lem:tail-symmetry-uniform} and \eqref{eq:swap-hellinger-identity} make
this quantity tend to zero as the tail recedes.  This contradiction proves
$\sum_{i>N_0}\operatorname{err}_i<\infty$.

Define $Q_0$ by the matched rows $\widehat e_i$ after the fixed initial cutoff
$N_0$, and by zero rows for $i\leq N_0$.
The summability just proved implies
\[
 \begin{aligned}
 \|T-Q_0\|_{\mathcal S_2}^2
 &=\sum_{i\leq N_0}\|T^*e_i\|_2^2
                      +\sum_{i>N_0}\|T^*e_i-\widehat e_i\|_2^2\\
 &\leq\sum_{i\leq N_0}\|T^*e_i\|_2^2
       +\sum_{i>N_0}\operatorname{err}_i<\infty.
 \end{aligned}
 \]
Thus $Q_0$ is a compact perturbation of the invertible operator $T$, hence
Fredholm of index zero.  Its cokernel is the span of the first $N_0$ output
coordinates and its kernel is spanned by the unused input coordinates, so
exactly $N_0$ columns are unused.  Matching them to the first $N_0$ rows with
positive signs completes $Q_0$ to $Q\in\mathcal P$ and proves
\eqref{eq:hilbert-schmidt-necessity}.

If $T$ is orthogonal, both covariance defects vanish, and every
$W=P T^*ST$ above is orthogonal.  Thus only the orthogonal part of
\cref{lem:local-hs-coercivity} is used, and it requires no scale Fisher
information.
\end{proof}

\subsection{Exact identification of the linear stabilizer}
\label{sec:exact-stabilizer}

Unlike measure-class equivalence, exact invariance requires only the second
moment.

\begin{proof}[Proof of \cref{cor:moment-only-stabilizer}]
Every member of $\mathcal P$ preserves the product law.
Conversely, if $\nu_T=\mu$, covariance gives $TT^*=I$.
Since $T$ is boundedly invertible, it is orthogonal.  Each of its unit
rows gives a linear combination with the same law as $X$.
Applying \cref{lem:unit-row-rigidity} to the constant sequence
consisting of this row shows that the row is an allowed signed unit
vector.  Orthogonality makes the selected coordinates a bijection.
Thus $T\in\mathcal P$.
\end{proof}
\subsection{Two-sided uniform integrability and completion of the proof}
\label{sec:likelihood-convergence}

The finite-coordinate limits and the two finite-marginal
conditional-expectation martingales now combine to complete the proof.

\begin{proposition}[Finite-coordinate approximants recover the likelihood ratio]
\label{prop:finite-coordinate-approximants-recover-likelihood-ratio}
Under the assumptions of \cref{thm:measure-equivalence}, suppose
\(\nu_T\sim\mu\).  For any \(Q\in\mathcal P\) satisfying
\(TQ^{-1}-I\in\mathcal S_2(H)\), put \(K=TQ^{-1}-I\) and
\(T_n=I+P_nKP_n\).  Then \(T_n\) is eventually invertible and
\begin{equation}
 \frac{d\nu_{T_n}}{d\mu}\longrightarrow
       \frac{d\nu_{TQ^{-1}}}{d\mu}
       \quad\text{in }L^1(\mu),
 \qquad
 \sqrt{\frac{d\nu_{T_n}}{d\mu}}
       \longrightarrow
 \sqrt{\frac{d\nu_{TQ^{-1}}}{d\mu}}
       \quad\text{in }L^2(\mu).
 \label{eq:likelihood-convergence}
\end{equation}
Since \(\nu_Q=\mu\), the limiting law is also \(\nu_T\).  Moreover,
the finite-marginal likelihood ratios and their reciprocals are uniformly
integrable in their respective reference measures.
\end{proposition}

\begin{proof}[Proof of \cref{thm:measure-equivalence} and \cref{prop:finite-coordinate-approximants-recover-likelihood-ratio}]
Necessity follows from \cref{lem:hilbert-schmidt-necessity}.  If
$T-Q\in\mathcal S_2$ for $Q\in\mathcal P$, then
$TQ^{-1}=I+(T-Q)Q^{-1}$ is invertible and differs from the identity
by a Hilbert--Schmidt operator.  By \cref{prop:hilbert-schmidt-sufficiency},
$\nu_{TQ^{-1}}\sim\mu$.  Since $Q_\#\mu=\mu$, this is
$\nu_T\sim\mu$.  The same proposition gives the finite-coordinate
likelihood-ratio and square-root limits stated in
\cref{prop:finite-coordinate-approximants-recover-likelihood-ratio}.
The exact linear stabilizer is identified in
\cref{cor:moment-only-stabilizer}.
Finally, \cref{lem:orbit-dichotomy,lem:finite-marginal-likelihood-ratios}
give the finite-marginal affinity characterization and the singular
alternative in \cref{thm:measure-equivalence}, and both
uniform-integrability statements in
\cref{prop:finite-coordinate-approximants-recover-likelihood-ratio}.
\end{proof}
\begin{proof}[Proof of \cref{cor:orthogonal-equivalence}]
Necessity is \cref{lem:hilbert-schmidt-necessity}.  For sufficiency, write
$UQ^{-1}=I+(U-Q)Q^{-1}$.  This is orthogonal and differs from the
identity by a Hilbert--Schmidt operator.  Apply
\cref{prop:orthogonal-sufficiency}, and then use
$Q_\#\mu=\mu$.  The alternative follows from
\cref{lem:orbit-dichotomy}.
\end{proof}

\subsection{Globalization of finite-dimensional coercivity}
\label{sec:globalization-proof}

Local coercivity controls only a fixed Frobenius neighborhood of the
stabilizer.  The global obstacle would be a sequence in varying dimensions
whose Hellinger displacement tends to zero while its distance from every
allowed signed permutation stays positive.  After passing to a subsequence
whose block Hellinger energies are summable, we stack these matrices as blocks
of one bounded invertible operator on \(\ell^2\).  Kakutani's theorem turns this
summability into equivalence of the infinite product, and the classification
forces the block-diagonal operator to be Hilbert--Schmidt-close to one allowed
signed permutation.  That
permutation must eventually preserve each block, contradicting the assumed
separation.  This direct-sum argument is the compactness input missing from the
local theorem.

\begin{proof}[Proof of \cref{thm:global-coercivity}]
Fix \(0<a\le b<\infty\).  We first prove uniform separation away from
the exact linear stabilizer: for every \(\varepsilon>0\),
\begin{equation}
 \inf\left\{
 h(\mu_n,A_\#\mu_n)^2:
 \begin{array}{l}
 n\ge1,\ A\in\GL(n,\R),\\[-2pt]
 s_{\min}(A)\ge a,\ s_{\max}(A)\le b,\\[-2pt]
 d_n(A)\ge\varepsilon
 \end{array}\right\}>0.
 \label{eq:uniform-separation}
\end{equation}
If this fails, choose matrices \(A_k\in\GL(n_k,\R)\) in the displayed class
whose Hellinger distances tend to zero.  Since
\[
 -2\log\Aff(\mu_{n_k},(A_k)_\#\mu_{n_k})
 =-2\log\left(1-\tfrac12
 h(\mu_{n_k},(A_k)_\#\mu_{n_k})^2\right),
 \]
we may pass to a subsequence for which the sum of these block energies is
finite.

On consecutive coordinate blocks define
$A=\bigoplus_{k\ge1}A_k\in\GL(H)$.  Each block law is equivalent to
$\mu_{n_k}$, and the energy summability gives
\[
 \prod_{k\ge1}\Aff\bigl(\mu_{n_k},(A_k)_\#\mu_{n_k}\bigr)>0.
\]
Kakutani's theorem \cite{kakutani1948equivalence} gives \(\nu_A\sim\mu\),
so \cref{thm:measure-equivalence} supplies \(Q\in\mathcal P\) with
\begin{equation}
 D:=A-Q\in\mathcal S_2(H).
 \label{eq:globalization-D}
\end{equation}

Let \(\Pi_k\) be the orthogonal projection onto the \(k\)-th output block.
The orthogonality of these blocks gives
\[
 \sum_k\|\Pi_kD\|_{\mathcal S_2}^2=\|D\|_{\mathcal S_2}^2<\infty,
 \qquad\text{hence}\qquad
 \|\Pi_kD\|_{\mathcal S_2}\longrightarrow0.
 \]
If $e_i$ belongs to the $k$th output block but $Q^*e_i$ lies outside the
corresponding input block, then the two vectors have disjoint supports and
\begin{equation}
 \|A^*e_i-Q^*e_i\|_2^2
 =\|A_k^*e_i\|_2^2+1\ge a^2+1.
 \label{eq:row-escape-cost}
 \end{equation}
For large $k$ this contradicts
$\|\Pi_kD\|_{\mathcal S_2}^2<a^2+1$.  Thus $Q^*$ preserves every sufficiently
large block; its restriction is a bijection and defines
$Q_k\in\mathcal P_{n_k}$.  Hence
\[
 d_{n_k}(A_k)
 \le\|A_k-Q_k\|_F
 =\|\Pi_kD\|_{\mathcal S_2}\longrightarrow0,
\]
contrary to $d_{n_k}(A_k)\ge\varepsilon$.  This proves
\eqref{eq:uniform-separation}.

Let $c_{\mathrm{loc}},C_{\mathrm{loc}},\delta_{\mathrm{loc}}$ be the local
constants, with $0<\delta_{\mathrm{loc}}<1$.  Choose a minimizer
$Q_n\in\mathcal P_n$ for $d_n(T)$.  Since
\((Q_n)_\#\mu_n=\mu_n\),
\[
 h(\mu_n,T_\#\mu_n)
 =h(\mu_n,(TQ_n^{-1})_\#\mu_n),
 \qquad
 \|TQ_n^{-1}-I\|_F=d_n(T).
\]
For $d_n(T)\le\delta_{\mathrm{loc}}$, the local theorem gives the required
bounds.  For $d_n(T)>\delta_{\mathrm{loc}}$,
\eqref{eq:uniform-separation} gives the lower bound and $h^2\le2$ the upper
bound.  Taking $c$ no larger than $c_{\mathrm{loc}}^2$ and the separation
constant at $\varepsilon=\delta_{\mathrm{loc}}$, and taking
\[
 C=\max\{C_{\mathrm{loc}}^2,2\delta_{\mathrm{loc}}^{-2}\}.
\]
proves \eqref{eq:global-coercivity}.
\end{proof}
